\documentclass[11pt,letterpaper]{article}
\usepackage[T1]{fontenc}
\usepackage{lmodern}
\usepackage[margin=0.94in,top=0.88in,bottom=0.87in,headheight=14pt]{geometry}
\usepackage{amsmath,amssymb,amsthm,booktabs,xcolor,fancyhdr}
\usepackage[colorlinks=true,urlcolor=teal,linkcolor=teal,citecolor=teal]{hyperref}
\definecolor{ink}{RGB}{22,53,72}
\definecolor{teal}{RGB}{5,94,98}
\hypersetup{pdftitle={Recovering all Dirichlet zeros from one zero-free interval},pdfauthor={OpenAI}}
\pagestyle{fancy}\fancyhf{}
\lhead{\small\color{ink}One-interval recovery of Dirichlet zeros}
\rhead{\small\color{ink}November 2026}
\cfoot{\small\thepage}
\renewcommand{\headrulewidth}{0.3pt}
\newtheorem{theorem}{Theorem}
\newtheorem{lemma}[theorem]{Lemma}
\newtheorem{corollary}[theorem]{Corollary}
\newcommand{\R}{\mathbb R}
\newcommand{\Z}{\mathbb Z}
\newcommand{\sech}{\operatorname{sech}}
\newcommand{\atanh}{\operatorname{artanh}}
\newcommand{\dd}{\mathrm d}
\newcommand{\ord}{\operatorname{ord}}
\setlength{\parindent}{1.15em}
\setlength{\parskip}{0.17em}
\setlength{\emergencystretch}{1em}
\numberwithin{equation}{section}
\begin{document}
\thispagestyle{plain}
\noindent{\sffamily\small\bfseries\color{teal} RESEARCH NOTE\hfill 26 NOVEMBER 2026}
\bigskip

\noindent{\LARGE\bfseries\color{ink} Recovering all Dirichlet zeros\\[3pt]from one zero-free interval\par}
\medskip
\noindent{\large An inverse theorem for the zeros of high logarithmic derivatives}
\bigskip

\noindent\textbf{Abstract.}
Choose a closed real interval containing no zero of a completed Dirichlet
$L$-function in the ordinate variable. Insert a double zero at its left
endpoint and a simple zero at its right endpoint. Each odd derivative of
order at least three of the logarithm of the absolute value then vanishes
at exactly one interior point. The locations of these points, even along
any unbounded subsequence of odd orders, determine every nontrivial zero
with multiplicity and determine the primitive character. For Riemann's
$\Xi$, one may use the fixed interval $[-1,1]$ without knowing any zero
in advance. A finite-order inequality bounds the next positive ordinate.
The proof also gives the exact exponential scale at which a changed
distant zero affects an observation.

\section{The observable}
Let $F(t)$ be the completed function on the critical line, for example
$F(t)=\Xi(t)=\xi(\tfrac12+it)$. Select $a<b$ with $F$ zero-free on
$[a,b]$. Form the entire function
\[
       G(t)=(t-a)^2(t-b)F(t).
\]
For every odd integer $p\ge3$, record the single real number $x_p\in(a,b)$
at which
\begin{equation}\label{obs}
       \frac{\dd^p}{\dd t^p}\log|G(t)|\bigg|_{t=x_p}=0.
\end{equation}
The data are the interval and the indexed point locations. They can be
obtained by differentiating the completed function within the interval.

\begin{theorem}[One-interval determination]\label{main}
Let $F$ be the entire completion of a primitive Dirichlet $L$-function in
the ordinate variable, including the Riemann completion in conductor one.
The point in \eqref{obs} exists and is unique. Given $a,b$ and the indexed
values $x_p$ for any unbounded set of odd $p\ge3$, the full multiset of
nontrivial zeros and the primitive character are uniquely determined.
\end{theorem}

The analytic input from Dirichlet GRH is that the zero ordinates are real.
The remaining argument concerns positive masses on a line. Two unequal
endpoint multiplicities break reflection symmetry. With equal endpoint
multiplicities, every spectrum symmetric about the midpoint would give
exactly the same midpoint at every order.

\begin{center}
\setlength{\unitlength}{0.93pt}
\begin{picture}(410,58)
\put(8,26){\vector(1,0){393}}
\color{teal}\linethickness{1.1pt}
\put(145,26){\line(1,0){120}}
\color{black}
\put(47,26){\circle*{5}}\put(47,12){\makebox(0,0){$c-R$}}
\put(145,26){\circle{10}}\put(145,26){\circle*{5}}
\put(145,10){\makebox(0,0){$a$}}
\put(205,23){\line(0,1){6}}\put(205,10){\makebox(0,0){$c$}}
\color{teal}\put(220,26){\circle*{5}}
\put(220,41){\makebox(0,0){$x_p$}}
\color{black}\put(265,26){\circle*{5}}\put(265,10){\makebox(0,0){$b$}}
\put(363,26){\circle*{5}}\put(363,12){\makebox(0,0){$c+R$}}
\put(47,42){\makebox(0,0){$L$}}\put(363,42){\makebox(0,0){$U$}}
\end{picture}\\[-1pt]
\small One observed position inside $(a,b)$ per order encodes the weights\\
at both exterior positions $c-R$ and $c+R$.
\end{center}
\newpage
\section{The real-variable inverse problem}
Write $a=c-d$, $b=c+d$, $d>0$. Let $\mu$ be a locally finite discrete
multiset of real locations $\gamma$ outside $[a,b]$, with positive weights
$m_\gamma$ and
\begin{equation}\label{sum}
            \sum_\gamma\frac{m_\gamma}{1+\gamma^2}<\infty.
\end{equation}
Here locally finite means finitely many distinct locations in each bounded
interval. Put known endpoint weights $A,B>0$ at $a,b$, and assume $A\ne B$.
For odd $p\ge3$ define
\begin{equation}\label{field}
 F_p(x)=\frac{A}{(x-a)^p}+\frac{B}{(x-b)^p}
                       +\sum_\gamma\frac{m_\gamma}{(x-\gamma)^p}.
\end{equation}
The sums and differentiated sums converge locally uniformly in $(a,b)$
by \eqref{sum}. Every term has negative derivative and the endpoint
limits are $+\infty$ and $-\infty$. Thus $F_p$ has exactly one zero $x_p$.
The zero with only the two endpoints present is
\begin{equation}\label{yp}
 y_p=c+d\tanh(h/p),\qquad h=\tfrac12\log(A/B).
\end{equation}
Indeed the ratio of the two endpoint distances at $y_p$ is $(A/B)^{1/p}$.

\begin{lemma}[Endpoint approximation]\label{endpoint}
If there is exterior mass and $R_0=\min_\gamma|\gamma-c|$, then
\[
        x_p=y_p+O\!\left(\frac{d}{p}(d/R_0)^p\right).
\]
If there is no exterior mass, $x_p=y_p$ for all $p$.
\end{lemma}
\begin{proof}
For any fixed $C$ and $S>d$, uniformly on $|x-c|\le C/p$,
\begin{equation}\label{tail}
 \sum_{|\gamma-c|\ge S}\frac{m_\gamma}{|x-\gamma|^p}=O(S^{-p}).
\end{equation}
For large $p$, replacing $|x-\gamma|^{-p}$ by $|c-\gamma|^{-p}$
costs at most $(1-C/(pS))^{-p}\le e^{2C/S}$. After extracting $S^{-p}$,
the remaining sum is bounded by $\sum m_\gamma(S/|c-\gamma|)^2$ over
the indicated locations. This is finite by \eqref{sum}. Apply \eqref{tail}
at $y_p$ with $S=R_0$. Throughout the gap, the nearer endpoint gives
$|F_p'(x)|\ge p\min(A,B)d^{-p-1}$. The mean value theorem proves the bound.
\end{proof}

The dimensionless displacement $u_p=(x_p-c)/d$ therefore agrees with
$\tanh(h/p)$ to exponentially small error. Its term of order $p^{-3}$
will distinguish a left-hand mass from a right-hand mass even when the
two contributions cancel to leading order. The requirement $h\ne0$
enters precisely at this step.
\newpage
\section{Recovering a radius and its two weights}
Suppose the locations at smaller radii are already known. Evaluate their
contribution and both endpoint terms at the observed $x_p$, and let $E_p$
be the negative of this known sum. Since $F_p(x_p)=0$, $E_p$ equals the
unknown exterior contribution. Let $R>d$ be its first occupied radius.
Write $L,U$ for the weights at $c-R,c+R$, assigning zero to an empty side.
If $S>R$ is the next occupied radius, put $\eta=R/S<1$. For a last radius
the tail terms below vanish. Define
\begin{equation}\label{polar}
 k=d/R,\qquad \theta_p=p\atanh(ku_p),\qquad
 G_p=(1-k^2u_p^2)^{-p/2},\qquad \theta_0=kh.
\end{equation}
Pairing the two locations at $R$ and applying \eqref{tail} gives
\begin{equation}\label{shell}
        \frac{R^pE_p}{G_p}
          =L e^{-\theta_p}-Ue^{\theta_p}+O(\eta^p).
\end{equation}
Indeed $(1+ku_p)^{-p}=G_p e^{-\theta_p}$ and $(1-ku_p)^{-p}=G_p e^{\theta_p}$.
Taylor expansion using Lemma~\ref{endpoint} yields
\begin{equation}\label{theta}
 \theta_p=kh+\frac{v}{p^2}+O(p^{-4}),\qquad
 v=\frac{(k^3-k)h^3}{3}\ne0,\qquad G_p=1+O(p^{-1}).
\end{equation}

\begin{theorem}[Explicit inversion]\label{inverse}
Every occupied radius and both of its weights are recovered by the following
successive limits. All limits may use any one unbounded subsequence of odd orders:
\begin{align}
 R&=\lim |E_p|^{-1/p},\qquad H=\lim Z_p,\qquad Z_p=R^p E_p/G_p,\label{root}\\
 J&=\lim\frac{Z_p-H}{\theta_p-\theta_0},\label{slope}\\
 L&=\tfrac12 e^{\theta_0}(H-J),\qquad
 U=-\tfrac12 e^{-\theta_0}(H+J).\label{masses}
\end{align}
\end{theorem}
\begin{proof}
Set $f(t)=L e^{-t}-Ue^t$. At $\theta_0$, $f=H$ and
$f'=-(L e^{-\theta_0}+Ue^{\theta_0})<0$. If $H\ne0$, the normalized
right side of \eqref{shell} tends to $H$. If $H=0$, it is
$f'(\theta_0)v/p^2+O(p^{-4})+O(\eta^p)$, with nonzero coefficient.
In either case $E_p$ is eventually nonzero and \eqref{root} recovers $R$.
Then \eqref{shell}--\eqref{theta} give $J=f'(\theta_0)$ and solving two
linear equations yields \eqref{masses}. Remove these weights and repeat.
Local finiteness puts only finitely many radii before any specified
location. An identically vanishing residual on the observed subsequence
characterizes an empty remainder by the same argument.
\end{proof}

These limits are statements about exact real observations. They impose no
lower density requirement on the orders. For example, retaining only
$p=2^n+1$, $n\ge1$, gives the same uniqueness conclusion. The next theorem
quantifies how close two observations can be even when the two spectra
are different.
\newpage
\section{How much a changed zero moves the observation}
Two measures $\mu,\widetilde\mu$ satisfying \eqref{sum} use the same
interval and the same unequal endpoint weights. Let $R$ be the nearest
radius at which they differ. Write $\Delta L,\Delta U$ for the weights of
$\mu-\widetilde\mu$ there, and use $k,\theta_0,v$ of
\eqref{polar}--\eqref{theta}. Set
\[
 H=\Delta L e^{-\theta_0}-\Delta U e^{\theta_0},\qquad
 J=-\Delta L e^{-\theta_0}-\Delta U e^{\theta_0}.
\]
\begin{theorem}[Resolution of the first discrepancy]\label{rate}
For the respective point sequences $x_p,\widetilde x_p$,
\[
 x_p-\widetilde x_p\sim
 \frac{d}{2\sqrt{AB}}\left(\frac dR\right)^p
 \begin{cases}
 H/p,&H\ne0,\\
 Jv/p^3,&H=0.
 \end{cases}
\]
In particular $\lim|x_p-\widetilde x_p|^{1/p}=d/R$.
\end{theorem}
\begin{proof}
Evaluate $F_p-\widetilde F_p$ at $\widetilde x_p$. Formula
\eqref{shell} applies to the signed difference by \eqref{tail}, with the
expansion for $\widetilde x_p$. If $H=0$ then $J\ne0$ since the two weight
differences are not both zero. The field difference is asymptotic to
$R^{-p}H$ or $R^{-p}Jv/p^2$. For any $w$ between the two observed points,
$p(w-c)/d\to h$ by Lemma~\ref{endpoint}, uniformly on this segment.
The two endpoint terms of $-d^{p+1} F'_p(w)/p$ tend respectively to
$Ae^{-h}$ and $Be^h$, whose sum is $2\sqrt{AB}$. All remaining terms tend
to zero by \eqref{tail} with exponent $p+1$. Apply the mean value theorem to
$F_p(x_p)-F_p(\widetilde x_p)$.
\end{proof}

\paragraph{An example where the leading contributions cancel.}
Use $(a,b)=(-1,1)$ and $(A,B)=(4,1)$. One measure has additional weights
$2,1$ at $-2,2$, the other has none. Then $H=0$, $J=-2\sqrt2$ and
$v=-(\log 2)^3/8$. With $y_p=\tanh((\log 2)/p)$,
\[
             p^3 2^p (x_p-y_p)\longrightarrow
             \frac{\sqrt2(\log 2)^3}{16}=0.02943549871\ldots .
\]
The included integer-arithmetic check clears positive denominators and
bisects both roots at dyadic rationals. Its certified enclosures are
\[
\begin{array}{c@{\qquad}l}
p&\text{enclosure of }p^3 2^p(x_p-y_p)\\ \hline
15 &[0.02900627349,\ 0.02900627350]\\
63 &[0.02934699772,\ 0.02934699773]\\
127 &[0.02939267493,\ 0.02939267494].
\end{array}
\]
The exponential factor quantifies the precision required for distant
zeros. The two different powers of $p$ quantify the additional loss caused
by cancellation within a pair.
\newpage
\section{Dirichlet zeros determine the character}
For a primitive character $\chi$ of conductor $q$, put
$e=\mathbf 1_{\{q=1\}}$ and $\epsilon=(1-\chi(-1))/2$. In the conductor-one
case $L=\zeta$ and $\epsilon=0$. Up to a constant, the entire completion is
\[
 C_\chi(s)=[s(s-1)]^e(q/\pi)^{(s+\epsilon)/2}
            \Gamma((s+\epsilon)/2)L(s,\chi).
\]
Its zeros, counted with multiplicity, are the nontrivial zeros. It has
order one. These classical properties follow from the theta integral and
functional equation; see Davenport \cite{davenport} and the DLMF
\cite{dlmf}. Dirichlet GRH places the zeros at
$\rho=\tfrac12+i\gamma$, $\gamma\in\R$. Order one implies
\eqref{sum}. Take $F(t)=C_\chi(\tfrac12+it)$. In its genus-at-most-one
Hadamard product the exponential factors have affine real logarithms on
the real line. Consequently, on the gap, for every odd $p\ge3$,
\[
 \frac{1}{(p-1)!}\frac{\dd^p}{\dd t^p}\log|G(t)|=F_p(t)
       \quad (A=2,\ B=1).
\]
This follows termwise by local uniform convergence after two
differentiations. A zero at $t=0$, when present, contributes its usual
power factor and the same formula. Theorem~\ref{inverse} recovers every
$\rho$, including multiplicity.

There is also a direct way to read the arithmetic data out of the recovered
zeros. For real $\sigma>1$, Hadamard factorization gives the absolutely
convergent sum
\[
 J_\chi(\sigma)=-\sum_\rho(\sigma-\rho)^{-2}
            =(\log C_\chi)''(\sigma).
\]
The derivatives here are with respect to $s$ and may be complex. Write
$\psi_1$ for the trigamma function. The Euler product gives
\[
 J_\chi(\sigma)=-e\big(\sigma^{-2}+(\sigma-1)^{-2}\big)
   +\tfrac14\psi_1((\sigma+\epsilon)/2)
   +\sum_{n\ge2} b_n n^{-\sigma},
 \qquad b_n=\Lambda(n)\chi(n)\log n ,
\]
where $\Lambda(n)$ is the von Mangoldt function. The series decays
exponentially as $\sigma\to+\infty$. From
$\psi_1(z)=z^{-1}+(2z^2)^{-1}+O(z^{-3})$ we obtain
\[
 \lim_{\sigma\to\infty}\sigma^2\left(J_\chi(\sigma)-\frac1{2\sigma}\right)
        =\frac{1-\epsilon}{2}-2e .
\]
The possible values $-3/2,1/2,0$ distinguish zeta, even nonprincipal,
and odd characters, respectively. Subtract the now known rational and
trigamma terms from $J_\chi$ to obtain $D(\sigma)=\sum b_n n^{-\sigma}$.
In increasing order of $n$ recover
\[
 b_n=\lim_{\sigma\to\infty} n^\sigma
       \left(D(\sigma)-\sum_{2\le j<n} b_j j^{-\sigma}\right).
\]
Absolute convergence at $\sigma=2$ justifies the limit by domination.
At a prime $\ell$, division by $(\log\ell)^2$ gives $\chi(\ell)$,
including zero values. Complete multiplicativity recovers $\chi$ on all
integers, and its unique primitive conductor follows. This proves
Theorem~\ref{main}.
\newpage
\section{Riemann's function on a fixed interval}
One can choose the two artificial zeros $-1,1$ before computing
any Riemann zero. To see the interval is zero-free, put
$P(v)=\sum_{n\ge1}e^{-\pi n^2v}$. The classical theta integral yields
\[
 \Xi(t)=\tfrac12-(t^2+\tfrac14)
      \int_1^\infty P(v)v^{-3/4}\cos((t/2)\log v)\,\dd v .
\]
For $v\ge1$, $P(v)\le e^{-\pi v}/(1-e^{-3\pi})$ since successive
squares differ by at least three. Hence on $[-1,1]$ the integral term
including its prefactor has absolute value at most
$5e^{-\pi}/[4\pi(1-e^{-3\pi})]<1/2$. In particular $\Xi(t)>0$ there.
Its zeros occur in pairs $\pm r$ of equal multiplicity $m_r$.

Let $h=(\log2)/2$. The field at zero equals $1$, and the sum from every
exterior pair is negative at $y_p=\tanh(h/p)$. Thus $0<x_p<y_p<h/p$.
After any number of the nearest pairs have been recovered, let $K_p$ be
the signed contribution of the two markers and the recovered pairs at
$x_p$. If $R$ is the next positive ordinate, then
\begin{equation}\label{positive}
 K_p=\sum_{r\ge R}m_r\big[(r-x_p)^{-p}-(r+x_p)^{-p}\big]>0.
\end{equation}
Every pair contributes positively to this readout.
Equations \eqref{shell}--\eqref{theta} simplify to
\begin{equation}\label{zetaform}
 R=\lim K_p^{-1/p},\qquad
 m_R=\frac{\lim R^pK_p}{2\sinh(h/R)} .
\end{equation}
For the first pair, $K_p=2/(1+x_p)^p-1/(1-x_p)^p$ uses only the
observed point. Each application of \eqref{zetaform} supplies the pair
required for the next application.

\subsection*{A finite two-order bound}
Define the positive decreasing kernel, for $r>z>0$,
\[
 A_j(r,z)=\frac{(r-z)^{-j}-(r+z)^{-j}}{2jz}
         =\tfrac12\int_{-1}^1(r+vz)^{-j-1}\,\dd v .
\]
Set $x=x_p$, $y=x_{p+2}$ and $D_j=K_j/(2j x_j)$, using the same recovered
pairs at both orders. Thus $D_j=\sum_{r\ge R}m_r A_j(r,x_j)$.
\begin{theorem}[Finite bounds]\label{bounds}
Let $\ell_p>x$ and $u_p^*\ge1$ solve the scalar equations
\[
 A_p(\ell_p,x)=D_p,\qquad
 Q_p(u_p^*)=D_{p+2}/D_p,\quad
 Q_p(r)=\frac{A_{p+2}(r,y)}{A_p(r,x)} .
\]
Both solutions exist and are unique in the stated domains. They satisfy
\[
     \ell_p\le R\le u_p^*,\qquad
           m_R\le D_p/A_p(u_p^*,x).
\]
\end{theorem}
\newpage
\begin{proof}
$A_p(r,x)$ decreases continuously from infinity to zero.
Since $m_R\ge1$, $D_p\ge A_p(R,x)$ and the first bound follows. From
the integral formula the logarithmic derivative of $A_j$ in $r$ lies
between $-(j+1)/(r-z)$ and $-(j+1)/(r+z)$. Hence
\[
 \frac{\dd}{\dd r}\log Q_p(r)
 \le-\frac{p+3}{r+y}+\frac{p+1}{r-x}<0\qquad(r\ge1).
\]
Indeed $(p+3)x+(p+1)y<3h<2\le2r$.
Now $D_{p+2}/D_p$ is a weighted average of $Q_p(r)$ over $r\ge R$
with positive weights $m_r A_p(r,x)$. It is at most $Q_p(R)$ and
strictly positive, while $Q_p(r)\to0$ at infinity. This proves existence,
uniqueness and the upper bound. Finally,
$D_p\ge m_R A_p(R,x)\ge m_R A_p(u_p^*,x)$.
\end{proof}

If $S>R$ is the next distinct positive ordinate, the upper bound satisfies
$u_p^*-R=O((R/S)^p)$. Indeed
\[
 D_j=m_R A_j(R,x_j)\{1+O((R/S)^j)\}
\]
by \eqref{tail} and \eqref{positive}, since $jx_j\to h$ and
$R^j A_j(R,x_j)\to\sinh(h/R)/h>0$. On $[R,R+1]$ the proof of
Theorem~\ref{bounds} bounds $-(\log Q_p)'$ below by a positive constant
uniformly in $p$. The log of the observed ratio thus differs from
$\log Q_p(R)$ by $O((R/S)^p)$, forcing the stated bound. If $m_R=1$,
the lower endpoint also satisfies $R-\ell_p=O((R/S)^p/p)$ using
$-\partial_r\log A_p\ge(p+1)/(r+x)$ between $\ell_p$ and $R$.
For a last pair both ratio measurements contain only that pair.
A multiplicity upper bound strictly below $2$ certifies that the first
remaining pair is simple.

\subsection*{Numerical illustration}
The program \texttt{xi\_demo.py} computes Taylor coefficients from complex
values of $\Xi$ by Cauchy's integral formula. No ordinates enter the
first experiment. All three quantities below use only the measurements
at orders $p,p+2$. The figures are rounded high-precision approximations.
\begin{center}\small
\begin{tabular}{rrrr}
\toprule
$p$ & $\ell_p$ & $u_p^*$ & upper bound on $m_R$\\ \midrule
 9 & 14.10262813980 & 14.22711060164 & 1.0918615321\\
17 & 14.13407780661 & 14.13795180923 & 1.0049451876\\
33 & 14.13472456790 & 14.13473049025 & 1.0000142459\\
65 & 14.1347251417338 & 14.1347251417509 & 1.00000000008\\
\bottomrule
\end{tabular}
\end{center}
The last interval has width $1.71\cdot10^{-11}$.
Its comparison value is $14.13472514173469\ldots$.
In a second experiment the known first pair is removed and its original
locations are used as the two endpoints. Orders 127 and 129 give the
numerical bounds $21.02203963874$ and $21.02203963945$ for the second
positive ordinate. The comparison value is $21.02203963877155\ldots$.
\newpage
\section{A reproducible interval certificate}
For orders 9 and 11 a separate program evaluates the observations by
interval arithmetic. It proves the following comparisons, with every
displayed endpoint rounded outwards:
\begin{align*}
0.038489153650346548707320 &\le x_9\le 0.038489153650346548707321,\\
0.031496268898746725992754 &\le x_{11}\le 0.031496268898746725992756,\\
1.09408109 &\le D_9/A_9(14.23,x_9)\le 1.09408110 .
\end{align*}
The program checks $D_9<A_9(14.10,x_9)$ and
$Q_9(14.23)<D_{11}/D_9$. Theorem~\ref{bounds} and the integrality of
multiplicity give the certificate
\[
          \boxed{14.10<\gamma_1<14.23,\qquad m_{\gamma_1}=1.}
\]
At orders 65 and 67 the same method gives the sharper certified result
\[
 \boxed{14.1347251417337<\gamma_1<14.1347251417510.}
\]
Here $x_{65}=0.0053318508623\ldots$ and $x_{67}=0.0051726940178\ldots$.
The computed upper bound on multiplicity is $1.000000000080332$.
The first zero and its simplicity are classical. The width
$1.73\cdot10^{-11}$ follows from two local measurements and the
proved bounds, with no zero ordinate supplied to the measurement routine.

\paragraph{How the certificate is checked.}
The substitution $v=\log u$ in the theta integral gives the kernel
$e^{v/4}P(e^v)\cos(tv/2)$. The program integrates its coefficients through
degree 44 on $0\le v\le5$, retaining $1\le n\le5$ in the theta sum.
It integrates the degree-19 Lagrange interpolant on panels of width $1/32$
using exact rational quadrature weights. On unit complex disks centered in
$[0,5]$, $\Re(e^v)>1/6$, so the analytic integrand for each coefficient,
including its factorial divisor, is bounded in modulus by $210$. Cauchy's
estimate and the interpolation remainder give an integrated error at most
$5\cdot210/32^{20}<10^{-27}$. The sum and integral truncations are also
bounded in the program. Its $10^{-26}$ enclosure for each coefficient
covers these three analytic errors. Integer arithmetic encloses all
algebraic operations, and the correctly rounded decimal exponential is
widened outwards.

On $|t|\le2$ the same integral gives $|\Xi(t)|<1$. This bounds the
discarded Taylor coefficients by $2^{-k}$ in degree $k$. The checker
shifts the enclosed series to each rational trial point, divides Taylor
series to evaluate the required coefficient of $\Xi'/\Xi$, and bisects
using signs certified by intervals. The resulting $x_9,x_{11}$
enclosures enter three rational comparisons. For orders 65 and 67 the program uses Taylor degree 182, 80 nodes per
panel, ten theta terms and 170 decimal places for interval endpoints.
The combined quadrature and truncation errors for each integrated
coefficient are below $10^{-116}$. The program checks signs of the full
field at both ends of each proposed bracket. Both modes are available in
\texttt{examples/xi\_certificate.py} with their recorded outputs.

\subsection*{Further research questions}
How should the marker weights and interval be chosen with prescribed
absolute observation errors? Can one obtain analogous finite bounds for
a complex character, whose positive and negative ordinates can carry
different masses? For zeta, reflection symmetry gives a positive
weighted average. For a general character, Theorem~\ref{inverse}
extracts the slope of a signed expression over an order-$p^{-2}$ change
in its argument.
\newpage
\section{Sources and prior work}
Sernau's Lean edition \cite{sernau} uses short forms, a spanning theorem
and a determinant bound to control a weighted prime-power remainder from
one side. A quadratic twist and Mellin residues connect that bound to
zero multiplicities. Its exported critical-strip theorem covers all
Dirichlet characters. We inspected its statements and reading guide.
Rubinstein independently reports a rebuild, kernel replay and comparison
of the compact edition against the five target statements \cite{rubinstein}.
The new inverse theorems have the written proofs above.

ProveIt's GRH index \cite{proveit} records 61 packets placed without an
independent mathematical review. The signed-zero-data packet recovers a
modulus and auxiliary prime from weighted ordinates. The prime-sampling
packet extracts multiplicity from weighted prime sums. The requested
encyclopedia article \cite{wiki} distinguishes Dirichlet GRH from other
critical-line conjectures. Only the Dirichlet statement enters here.

P\'olya, Shaw--Prather and B\o gvad--H\"agg \cite{polya,shaw,bogvad}
study where zeros of high derivatives accumulate for a given singularity
configuration. Abel--Goncharov interpolation \cite{popov} studies
uniqueness from derivative data for entire functions. Kawalec
\cite{kawalec} recovers zeta ordinates recursively from secondary zeta
values, and Matiyasevich \cite{mati} studies recovery from local derivative
values. The proposed contribution here is recovery from point locations
at arbitrarily sparse odd orders, including the cancellation law and
finite two-order bounds. A targeted search and the inspected packets did
not identify these statements. Historical priority remains open.

\small
\begin{thebibliography}{10}\setlength{\itemsep}{1pt}\setlength{\parsep}{0pt}\setlength{\parskip}{0pt}
\bibitem{sernau} L. Sernau et al., \href{https://github.com/lsernau/grh-lean/tree/72fead102f663d5c7a407964a13dc0cf9d977207}{\emph{GRH for Dirichlet $L$-functions, in Lean 4}},
commit \texttt{72fead1}. See \texttt{grh/GUIDE.md}, \texttt{grh/src/Statements.lean}, \texttt{grh/src/Main.lean} and \texttt{New/Reduction/Zeros/Criterion.lean} under \texttt{grh/src/}.
\bibitem{rubinstein} M. Rubinstein, \href{https://www.math.uwaterloo.ca/~mrubinst/alpoge_claude_grh/}{\emph{Alp\"oge/Claude proof of the Generalized Riemann Hypothesis}}, verification report, 26 Nov. 2026.
\bibitem{proveit} V. Reshetnikov, \href{https://github.com/VladimirReshetnikov/ProveIt/tree/a9ae674322c8c4cce1bb2262e577733d3959d741/Analysis/GeneralizedRiemannHypothesis}{\emph{ProveIt, Analysis/GeneralizedRiemannHypothesis}}, commit \texttt{a9ae674}. See \href{https://github.com/VladimirReshetnikov/ProveIt/blob/a9ae674322c8c4cce1bb2262e577733d3959d741/Analysis/GeneralizedRiemannHypothesis/Research/grh-torsion-and-zero-data/recovery.tex}{the signed-zero reconstruction} and \href{https://github.com/VladimirReshetnikov/ProveIt/blob/a9ae674322c8c4cce1bb2262e577733d3959d741/Analysis/GeneralizedRiemannHypothesis/Research/grh-torsion-and-prime-sampling/prime-sampling.tex}{the prime-sampling packet}.
\bibitem{wiki} Wikipedia contributors, \href{https://en.wikipedia.org/wiki/Riemann_hypothesis}{\emph{Riemann hypothesis}}, consulted 26 Nov. 2026.
\bibitem{davenport} H. Davenport, \emph{Multiplicative Number Theory}, 3rd ed., revised by H. L. Montgomery, Springer, 2000, chs. 9--12.
\bibitem{dlmf} NIST, \emph{Digital Library of Mathematical Functions}, \href{https://dlmf.nist.gov/25.4}{\S25.4} and \href{https://dlmf.nist.gov/25.15}{\S25.15}.
\bibitem{polya} G. P\'olya, \href{https://projecteuclid.org/journals/bulletin-of-the-american-mathematical-society/volume-49/issue-3.P1/On-the-zeros-of-the-derivatives-of-a-function-and/bams/1183505036.full}{On the zeros of the derivatives of a function and its analytic character}, \emph{Bull. Amer. Math. Soc.} 49 (1943), 178--191.
\bibitem{shaw} J. K. Shaw and C. L. Prather, \href{https://doi.org/10.1155/S0161171282000635}{A P\'olya shire theorem for functions with algebraic singularities}, \emph{Int. J. Math. Math. Sci.} 5 (1982), 691--706.
\bibitem{bogvad} R. B\o gvad and C. H\"agg, \href{https://arxiv.org/abs/1610.00921}{A refinement for rational functions of P\'olya's method to construct Voronoi diagrams}, \emph{J. Math. Anal. Appl.} 452 (2017), 312--334.
\bibitem{kawalec} A. Kawalec, \href{https://arxiv.org/abs/2012.06581}{Analytical recurrence formulas for non-trivial zeros of the Riemann zeta function}, arXiv:2012.06581v5 (2022).
\bibitem{mati} Yu. Matiyasevich, \href{https://doi.org/10.1016/j.ic.2023.105130}{Towards non-iterative calculation of the zeros of the Riemann zeta function}, \emph{Information and Computation} 296 (2024), 105130.
\bibitem{popov} A. Yu. Popov, \href{https://geodesic.mathdoc.fr/item/SM_2002_193_2_a4/}{Bounds for convergence and uniqueness in Abel--Goncharov interpolation}, \emph{Sb. Math.} 193 (2002), 247--277.
\end{thebibliography}
\end{document}
